\documentclass[11pt]{ctexart}
\usepackage[a4paper,margin=1in]{geometry}
\usepackage{graphicx}
\usepackage{xcolor}
\usepackage{hyperref}
\definecolor{refgray}{gray}{0.45}
\title{\textbf{Market Views｜全球投行叙事汇编}\\[0.4em]{\large 中国通胀平台期与全球利率的能源扰动，AI基础设施需求结构分化加剧}}
\author{KC桌面 自动生成}
\date{260712}
\begin{document}
\maketitle
\tableofcontents
\newpage
\section{一页摘要}
\begin{itemize}
  \item 中国6月CPI和PPI环比均回落，核心CPI创9个月新低，再通胀进程暂停而非逆转，能源价格回落是主要拖累。
  \item 全球利率市场主导变量从AI借贷叙事切换回能源价格，美债长端估值回归合理，近端面临CPI数据前的非对称风险。
  \item 人民币在美元走强背景下维持渐进升值路径，G10利差交易处于高利差低波动黄金窗口。
  \item AI基础设施投资持续高景气，但需求结构从全面扩张转向高度集中，资金与订单集中于数据层、AI网络与存储，传统软件承压。
  \item 电力供应成为数据中心扩张核心瓶颈，推动燃气轮机、固态变压器等新品类需求爆发。
  \item 中国运动用品需求未见结构性改善，批发渠道数据持续疲弱，宝胜6月零售额同比下滑8.6\%。
  \item 中国新能源乘用车渗透率6月突破62\%，但内需筑底，出口成为核心增长支柱，欧洲取代亚洲成为第一大出口目的地。
  \item 中国商业航天在长征十号B网捕回收成功后进入成本拐点，但发射成本仍远高于猎鹰9号。
\end{itemize}
\section{全球投行叙事汇编}
今日纳入 40 篇投行/券商研究，按 7 个市场、资产与产业主题系统整理。
\newpage
\subsection{宏观与利率：中国通胀平台期与全球利率的能源扰动}
\textbf{中国再通胀进程进入平台期，核心CPI创9个月新低，但AI相关领域和猪周期底部提供零星支撑；全球利率市场的主导变量从AI借贷叙事切换回能源价格，美债长端估值已回归合理，近端利率面临CPI数据前的非对称风险。}
\subsubsection*{主流共识}
\begin{itemize}
  \item 中国6月CPI和PPI环比均回落，再通胀进程暂停而非逆转，国内需求走软对价格动能形成实质性侵蚀。
  \item 能源价格回落是6月通胀低于预期的主要拖累，中东冲突溢价快速回吐，石油相关PPI环比骤降11.8\%。
  \item 核心CPI同比降至1.0\%，为去年9月以来最低，环比连续第二个月下滑，剔除能源和食品后的价格韧性减弱。
\end{itemize}
\subsubsection*{分歧与条件差异}
\begin{itemize}
  \item 财政紧缩的来源：Citi认为中央财政紧缩来自收入超预期增长（个税+12.2\%，印花税+88.9\%），而非支出削减；地方则受收入下滑和支出约束双重挤压。
  \item AI借贷对美债收益率的推升作用：GS认为市场高估了AI借贷的影响，长端收益率与宏观基本面之间并无明显脱节；而部分市场参与者仍将其视为关键驱动。
  \item 通胀前景：DB认为补贴效应递减将令消费品价格进一步承压；NOM和Citi则关注AI相关领域（通信设备CPI创历史新高7.6\%）和猪周期底部提供的结构性支撑。
\end{itemize}
\subsubsection*{机构观点}
\paragraph{Citi} 二季度宏观环境节奏变化超预期，中央与地方财政调整机制不同：中央因收入超预期（个税+12.2\%，印花税+88.9\%）形成事实紧缩，广义赤字率从2.4\%降至2.2\%；地方受土地出让收入同比-28.7\%和支出意愿审慎双重约束，支出同比-1.7\%。下半年财政部署加速将自然逆转这一调整。通胀方面，6月CPI同比1.0\%低于预期，能源冲击消退，但AI相关领域（电子制造PPI同比3.3\%创新高）和猪周期底部提供零星支撑，二季度名义GDP增速可能继续上行。
{\small\color{refgray}数据：前5个月中央广义支出同比+10.1\%，进度五年同期最高\par}
{\small\color{refgray}数据：土地出让收入12个月滚动值3.8万亿元，较2021年峰值9万亿元大幅回落\par}
{\small\color{refgray}数据：6月CPI同比1.0\%，PPI同比4.1\%，核心CPI同比1.0\%\par}
{\small\color{refgray}数据：电子制造PPI同比3.3\%，创1996年以来新高\par}
{\small\color{refgray}边际变化：市场此前将财政紧缩等同于支出削减，但Citi指出中央紧缩来自收入端超预期，且下半年将自然逆转；通胀叙事从总量关注转向结构分化，AI相关领域成为新的价格韧性来源。\par}
\paragraph{DB} 中国再通胀进程进入‘中场休息’，6月CPI环比-0.3\%，核心CPI同比降至1.0\%（去年9月以来最低），环比连续第二个月下滑。消费品价格走弱，家电和汽车价格受以旧换新补贴边际效应递减影响，同比涨幅回落0.5个百分点至1.1\%。PPI同比虽因基数升至4.1\%，但环比转负至-0.3\%，下游价格传导有限，反映需求不足。服务通胀持稳于0.8\%，通信设备价格同比升至7.6\%是少数韧性点。预计全年CPI 1.3\%、PPI 3.0\%，下半年需更多政策支持。
{\small\color{refgray}数据：6月CPI环比-0.3\%，核心CPI同比1.0\%创9个月新低\par}
{\small\color{refgray}数据：消费品价格同比涨幅回落0.5ppt至1.1\%\par}
{\small\color{refgray}数据：PPI环比-0.3\%，同比4.1\%\par}
{\small\color{refgray}数据：通信设备CPI同比7.6\%，较上月再提高1.0ppt\par}
{\small\color{refgray}边际变化：此前市场关注能源价格对通胀的拉动，DB指出补贴效应递减正成为消费品价格的新下行压力，而AI相关产品价格是少数上行亮点。\par}
\paragraph{GS} 美债长端估值已从‘略贵’回归至‘合理’，但AI相关企业借贷推升收益率的逻辑被高估——收益率水平与宏观及财政基本面之间并无明显脱节。当前全球利率市场的核心驱动变量是能源价格，美伊冲突升级使油价与美债收益率几乎同步波动。近端利率面临CPI数据前的非对称风险：油价越高，市场对偏强通胀数据的脆弱性越大。欧洲和英国前端利率的抛售提醒反通胀停滞对多头仓位的风险。
{\small\color{refgray}数据：市场定价到2027年初约45bp加息\par}
{\small\color{refgray}数据：美债长端估值从‘略贵’修复至‘大致公允’\par}
{\small\color{refgray}数据：欧洲前端利率因能源价格涨势戛然而止\par}
{\small\color{refgray}数据：英国曲线陡峭化交易优于直接做空Gilts\par}
{\small\color{refgray}边际变化：市场此前聚焦AI借贷作为美债收益率上行的主因，GS指出能源价格已重新成为主导变量，且AI借贷的影响被高估。\par}
\paragraph{NOM} 6月CPI同比1.0\%低于预期，环比连续第二个月下降，消费端疲软持续，但政策端已开始结构性回应——商务部等九部门发布零售业创新发展意见，目标2030年建立现代零售体系，包括土地、融资、减税支持及支持优质新零售企业上市。PPI同比升至4.1\%，‘上游热、下游冷’格局延续。消费回暖需时间，机构偏好战略清晰、利润率可见的领域，如运动服饰（安踏约14倍PE）和黄金零售（老铺黄金约7.7倍PE）。
{\small\color{refgray}数据：6月CPI同比1.0\%，环比-0.3\%连续第二个月下降\par}
{\small\color{refgray}数据：PPI同比4.1\%，环比转负\par}
{\small\color{refgray}数据：九部门联合发文支持零售业，目标2030年现代零售体系\par}
{\small\color{refgray}数据：安踏估值约14倍PE，老铺黄金约7.7倍PE\par}
{\small\color{refgray}边际变化：此前市场仅关注CPI读数走弱，NOM强调政策端已出现结构性回应，且消费领域存在估值有吸引力的结构性机会。\par}
\subsubsection*{关键数据}
\begin{itemize}
  \item 中国6月CPI同比1.0\%，核心CPI同比1.0\%创9个月新低
  \item PPI同比4.1\%，环比-0.3\%
  \item 中央广义支出前5个月同比+10.1\%，进度五年最高
  \item 土地出让收入12个月滚动值3.8万亿元，较峰值9万亿大幅回落
  \item 通信设备CPI同比7.6\%创历史新高
  \item 电子制造PPI同比3.3\%创1996年以来新高
  \item 美债长端估值从‘略贵’修复至‘大致公允’
  \item 市场定价到2027年初约45bp加息
\end{itemize}
\begin{figure}[htbp]
\centering
\includegraphics[width=0.88\linewidth]{figures/fig_017.jpg}
\caption{Figure 1 - ■ Central government austerity is a less obvious but a meaningful contributor – driven by revenue outperformance rather than spending restraint. Broad central government spending is running ahead of the run-rate seen in recent years, reaching a five-year high }
\end{figure}
\begin{figure}[htbp]
\centering
\includegraphics[width=0.88\linewidth]{figures/fig_022.jpg}
\caption{Figure 1 - ■ Domestic energy prices: Coal prices rose into the summer and partially offset the drag from global prices. PPI for coal extraction inched up to 5.6\%MoM with its year-on-year change at 20.6\%YoY after turning positive only three months ago. Coal prices could h}
\end{figure}
\begin{figure}[htbp]
\centering
\includegraphics[width=0.88\linewidth]{figures/fig_039.jpg}
\caption{Figure 1 - Figure 1: Momentum slowed across upstream and midstream Figure 2: Limited pass-through to downstream industries emerged Figure 3: Energy prices' contribution to CPI dropped by 0.2ppt Jan-24 Apr-24 Jul-24 Oct-24 Jan-25 Apr-25}
\end{figure}
\begin{figure}[htbp]
\centering
\includegraphics[width=0.88\linewidth]{figures/fig_060.jpg}
\caption{Exhibit 2 - Exhibit 1: US yields traded with oil prices for much of the week Exhibit 2: Term bills have cheapened vs OIS relative to shorter term bills as MMFs skewed their WAM shorter Dollar strength still a drag on foreign official deman}
\end{figure}
{\small\color{refgray}本节综合 6 篇研究；涉及 Citi、DB、GS、NOM。}
\newpage
\subsection{外汇与大宗商品：分化与错位中的结构性机会}
\textbf{当前外汇与大宗商品市场呈现显著分化：人民币在美元走强背景下维持升值路径，G10利差交易处于高利差低波动黄金窗口；锂价下跌但库存反降，市场误判需求韧性；金价回调暴露固定定价模式脆弱性，但铜金双主线逻辑未变。}
\subsubsection*{主流共识}
\begin{itemize}
  \item 人民币在出口强劲和政策支持下，即使美元走强仍维持渐进升值路径。
  \item G10外汇市场当前高利差与低波动组合使利差交易策略吸引力达近二十年极值。
  \item 锂价短期承压，但库存和产量下降显示边际改善，情绪驱动大于基本面转折。
\end{itemize}
\subsubsection*{分歧与条件差异}
\begin{itemize}
  \item 市场将锂矿复产视为供给冲击，但Citi认为复产更可能是需求验证信号。
  \item 金价回调对固定定价模式（如老铺黄金）的销售冲击与毛利率改善并存，市场对品牌定价权韧性存在分歧。
  \item 紫金矿业上半年产量完成率不足50\%，但JPM认为下半年集中释放可追赶指引，市场对产量节奏判断存在分歧。
  \item 美元走强下，人民币与欧元、新加坡元等低息货币表现分化，GS预计人民币将跑赢。
\end{itemize}
\subsubsection*{机构观点}
\paragraph{Citi} Citi认为市场对锂价下跌反应过度，两个触发因素（江西钨业复产、LFP正极产能调控提案）实际影响有限。碳酸锂现货从16.25万元/吨降至15.85万元/吨，但上周库存环比下降2\%至9.22万吨，冶炼厂库存降9\%，产量降3\%。宁德时代业绩会可能提供正面信号，复产细节和出货指引将验证需求韧性。
{\small\color{refgray}数据：碳酸锂现货从16.25万元降至15.85万元/吨\par}
{\small\color{refgray}数据：氢氧化锂从15.85万元降至14.25万元/吨\par}
{\small\color{refgray}数据：上周碳酸锂产量2.49万吨，环比-3\%\par}
{\small\color{refgray}数据：总库存9.22万吨，环比-2\%\par}
{\small\color{refgray}边际变化：市场此前聚焦复产和产能调控提案的供给冲击，Citi转向强调需求验证和库存下降的边际改善。\par}
\paragraph{GS} GS认为人民币升值路径未因美元走强而中断，维持12个月美元兑人民币6.50目标，预计人民币跑赢欧元和新加坡元。中间价在美元偏强背景下仍创2023年初以来新低至6.80，显示政策接受度提高。G10外汇市场处于高利差低波动窗口，利差交易策略吸引力达近二十年极值，但美元内部分化明显——对欧元日元走强，对部分新兴市场未必。
{\small\color{refgray}数据：人民币中间价自2023年初以来首次跌破6.80\par}
{\small\color{refgray}数据：GS 12个月美元兑人民币目标6.50\par}
{\small\color{refgray}数据：G10利差交易吸引力达近二十年极值\par}
{\small\color{refgray}数据：美元对欧元和日元走强，对部分新兴市场货币未必\par}
{\small\color{refgray}边际变化：此前市场认为美元走强将中断人民币升值，GS强调政策接受度提高和出口数据支撑，同时指出G10利差交易窗口罕见。\par}
\paragraph{GS} GS指出金价从5000美元回调20\%至4100美元，暴露老铺黄金固定定价模式脆弱性：天猫旗舰店二季度销售额同比下滑63\%（一季度+155\%），但毛利率反而受益，若金价维持4000美元，折扣后毛利率可达52\%远超40\%目标。品牌热度健康，公司已通过新品定价和促销测试更低价格。GS将2026年盈利预测下调9\%至80亿元，但认为当前7倍市盈率已反映不确定性。
{\small\color{refgray}数据：COMEX黄金从5000美元回调至4100美元，跌幅超20\%\par}
{\small\color{refgray}数据：老铺黄金天猫旗舰店二季度销售额同比-63\%（一季度+155\%）\par}
{\small\color{refgray}数据：2月提价后每克溢价从14\%-44\%扩大至81\%-131\%\par}
{\small\color{refgray}数据：若金价4000美元，折扣后毛利率可达52\%\par}
{\small\color{refgray}边际变化：此前市场关注金价上涨时老铺黄金的利润弹性，GS转向分析金价回调对固定定价模式的销售冲击与毛利率改善并存。\par}
\paragraph{JPM} JPM认为紫金矿业二季度核心净利润194亿元（环比+5\%，同比+66\%）低于预期215亿元，主因金价环比跌9\%、产量完成率不足50\%及成本上升。上半年黄金47吨（完成全年指引42\%）、铜53.4万吨（48\%）、锂4.3万吨（35\%）。但JPM仍列为首选，因产量集中在下半年释放，Manono锂矿项目预计下半年贡献更明显，成本信息待管理层更新。
{\small\color{refgray}数据：紫金矿业二季度核心净利润194亿元，低于JPM预期215亿元\par}
{\small\color{refgray}数据：环比+5\%，同比+66\%\par}
{\small\color{refgray}数据：上半年黄金产量47吨，完成全年指引42\%\par}
{\small\color{refgray}数据：铜产量53.4万吨，完成48\%\par}
{\small\color{refgray}边际变化：此前市场预期紫金矿业上半年产量进度正常，JPM指出完成率不足50\%但强调下半年集中释放，成本上升成为新关注点。\par}
\subsubsection*{关键数据}
\begin{itemize}
  \item 碳酸锂现货从16.25万元降至15.85万元/吨，库存环比-2\%
  \item 人民币中间价自2023年初以来首次跌破6.80
  \item G10利差交易吸引力达近二十年极值
  \item COMEX黄金从5000美元回调至4100美元，跌幅超20\%
  \item 老铺黄金天猫旗舰店二季度销售额同比-63\%
  \item 紫金矿业二季度核心净利润194亿元，低于预期215亿元
  \item 紫金矿业上半年黄金产量完成全年指引42\%
  \item 紫金黄金二季度盈利5.9亿美元，低于预期8.3亿美元
\end{itemize}
\begin{figure}[htbp]
\centering
\includegraphics[width=0.88\linewidth]{figures/fig_012.jpg}
\caption{Figure 1 - \#\# Shreyas Madabushi Figure 1. Lithium carbonate monthly inventory © 2026 Citi Inc. No redistribution without Citi's written permission. Figure 2. Lithium carbonate weekly inventory © 2026 Citi Inc. No redistribution without Citi's written permission.}
\end{figure}
\begin{figure}[htbp]
\centering
\includegraphics[width=0.88\linewidth]{figures/fig_057.jpg}
\caption{Exhibit 3 - Exhibit 4: G10 FX markets currently boast the combination of historically elevated carry levels with historically subdued vol...}
\end{figure}
\begin{figure}[htbp]
\centering
\includegraphics[width=0.88\linewidth]{figures/fig_065.jpg}
\caption{Exhibit 1 - Exhibit 1: Laopu historical price hike Exhibit 2: GS expects COMEX gold price to reach \textbackslash{}\$4,900 by Dec-26, implying \textbackslash{}\textasciitilde{}18\% upside from now Exhibit 3: Post the Feb price hike, Laopu is able to reach 40\% GPM when gold price at \textbackslash{}\textasciitilde{}US}
\end{figure}
\begin{figure}[htbp]
\centering
\includegraphics[width=0.88\linewidth]{figures/fig_013.jpg}
\caption{Figure 1 - Figure 1. Lithium carbonate monthly inventory © 2026 Citi Inc. No redistribution without Citi's written permission. Figure 2. Lithium carbonate weekly inventory © 2026 Citi Inc. No redistribution without Citi's written permission. Figure 3. Spread of battery-g}
\end{figure}
{\small\color{refgray}本节综合 4 篇研究；涉及 Citi、GS、JPM。}
\newpage
\subsection{AI基础设施：需求结构分化与电力瓶颈凸显}
\textbf{AI基础设施投资持续高景气，但需求结构正从‘全面扩张’转向‘高度集中’：资金与订单集中于数据层、AI网络与存储，传统软件与部分硬件承压；电力供应成为数据中心扩张的核心约束，推动燃气轮机、固态变压器等新品类需求爆发。}
\subsubsection*{主流共识}
\begin{itemize}
  \item AI基础设施投资持续增长，数据中心仍是AI落地的关键载体，IT预算增速回升至+3.3\%。
  \item AI支出高度集中于数据基础设施与AI网络，传统应用软件面临预算挤占与ROI审视。
  \item 电力供应成为数据中心扩张的核心瓶颈，推动燃气轮机、固态变压器等新品类需求增长。
\end{itemize}
\subsubsection*{分歧与条件差异}
\begin{itemize}
  \item AI需求是否见顶：部分市场担忧微软取消数据中心协议，但GS专家认为电力约束而非需求见顶，Agentic AI与企业内部数据训练提供双重支撑。
  \item AI支出来源：新增预算占比下降（从73\%降至69\%），挤占效应上升（50\%的CIO认为AI支出影响其他IT预算），市场对传统软件复苏节奏存在分歧。
  \item 存储超级周期持续性：GS认为记忆体超级周期支撑韩国与台湾市场资金流入，但伴随高波动性；Bernstein认为AI ASIC增长前景被低估，联发科等公司估值体系需重估。
  \item 太空AI商业化节奏：Citi认为工程问题已解决，但成本仍是瓶颈；发射能力扩张需监管与供应链同步跟进。
\end{itemize}
\subsubsection*{机构观点}
\paragraph{Bernstein} 大立光CPO光纤阵列组件进展超前，自动化试产线最早2027年中贡献收入，市场估值框架需从‘手机镜头周期股’切换至‘光学平台+CPO新品类’。手机镜头基本盘仍在升级，iPhone 18可变光圈镜头3Q量产，是近年最具挑战性的高端换代。毛利率49.4\%略低于预期，属新品爬坡正常波动。
{\small\color{refgray}数据：CPO自动化试产线最早2027年中贡献收入\par}
{\small\color{refgray}数据：2Q26毛利率49.4\%，略低于预期0.4个百分点\par}
{\small\color{refgray}数据：iPhone 18可变光圈镜头3Q量产，颗粒控制要求比潜望镜更严格\par}
{\small\color{refgray}数据：CPO FA组件价值量可能高于镜头\par}
{\small\color{refgray}边际变化：CPO业务进展从‘概念验证’进入‘样品交付与试产线建设’阶段，市场对估值框架的切换预期提前。\par}
\paragraph{Bernstein} 联发科2Q26营收1520亿新台币，超预期约5\%，但结构切换更关键：AI ASIC增长已足以覆盖手机端调整，竞争环境比预期更有利。手机业务面临内存成本上升与提价10-15\%带来的需求调整，但AI ASIC 2026-2028年收入预计达20亿/150亿/220亿美元，TPU部署范围扩大。报价回调15\%属正常调整，核心逻辑仍在AI ASIC。
{\small\color{refgray}数据：2Q26营收NT\$1520亿，超预期约5\%\par}
{\small\color{refgray}数据：AI ASIC 2026-2028年收入预计20亿/150亿/220亿美元\par}
{\small\color{refgray}数据：手机SoC提价10-15\%应对成本上升\par}
{\small\color{refgray}数据：报价从5月底高点回调约15\%\par}
{\small\color{refgray}边际变化：AI ASIC增长从‘补充’变为‘主导’，估值体系需从手机芯片公司向AI芯片公司切换。\par}
\paragraph{Citi} 东方电气燃气轮机正从数据中心需求中获取超额订单，2026年上半年签署10台50MW燃气轮机出口加拿大数据中心，单价1.5-2亿元，毛利率超30\%。产能10-12台/年，一期扩产至30台/年（2027年底），在手订单20台已覆盖未来两年产能。但市场对出口收入预期可能过于乐观，2026-2027年每年约20亿元，仅占总营收不到3\%。
{\small\color{refgray}数据：2026年上半年签署10台50MW燃气轮机出口订单\par}
{\small\color{refgray}数据：单台价格1.5-2亿元，毛利率超30\%\par}
{\small\color{refgray}数据：当前产能10-12台/年，扩产后2027年底达30台/年\par}
{\small\color{refgray}数据：在手订单20台，交付周期12个月\par}
{\small\color{refgray}边际变化：数据中心对燃气轮机的需求从‘概念’转向‘订单落地’，加拿大订单已签，美国市场是下一个变量。\par}
\paragraph{Citi} 2Q26 CIO调查显示全球IT预算增速回升至+3.3\%，但AI支出正从‘额外投入’变成‘替代性投入’：50\%的CIO认为AI支出挤占其他IT预算（3月为40\%）。AI预算占比6.5\%，未来12个月再增10\%，但新增预算占比从73\%降至69\%。数据与AI蝉联优先级第一，网络安全第二，Web安全首次进入高质量优先级。
{\small\color{refgray}数据：全球IT预算未来12个月增速+3.3\%（上期+2.6\%）\par}
{\small\color{refgray}数据：AI预算占比6.5\%，未来12个月再增10\%\par}
{\small\color{refgray}数据：50\%的CIO认为AI支出挤占其他IT预算（3月40\%，12月30\%）\par}
{\small\color{refgray}数据：新增预算占比从73\%降至69\%\par}
{\small\color{refgray}边际变化：AI支出从‘增量叠加’转向‘存量替代’，对传统软件公司的结构性挑战加剧。\par}
\paragraph{Citi} 大族激光与大族数控2Q26业绩远超预期，AI服务器需求是核心引擎。大族激光净利润同比增长176\%-207\%，IT设备业务（苹果供应链）1H26同比增长180\%至23亿元；大族数控净利润同比增长294\%-362\%，AI PCB设备（CCD背钻、激光钻孔）是主要推动力。AI驱动的PCB设备需求周期已从概念验证进入业绩兑现阶段。
{\small\color{refgray}数据：大族激光2Q26净利润同比增长176\%-207\%\par}
{\small\color{refgray}数据：大族数控2Q26净利润同比增长294\%-362\%\par}
{\small\color{refgray}数据：大族激光IT设备业务1H26同比增长180\%至23亿元\par}
{\small\color{refgray}数据：大族数控1H26营收同比至少增长至48亿元\par}
{\small\color{refgray}边际变化：AI PCB设备需求从‘概念验证’进入‘业绩兑现阶段’，订单与营收同步高增。\par}
\paragraph{Citi} 太空AI工程问题已解决，瓶颈在于成本。轨道AI卫星关键工程问题已解决，但先进材料和组件成本过高。发射能力有望大幅扩张，前FAA官员认为每年数千次发射可行，但需监管与供应链同步跟进。太空市场正进入10年以上研究周期，类似互联网时代，政府和商业投入相互强化。
{\small\color{refgray}数据：轨道AI卫星工程问题已解决，成本是主要瓶颈\par}
{\small\color{refgray}数据：前FAA官员认为每年数千次发射可行\par}
{\small\color{refgray}数据：太空市场研究周期可能持续10年以上\par}
{\small\color{refgray}数据：开源模型已达18个月前前沿水平\par}
{\small\color{refgray}边际变化：太空AI从‘愿景驱动’转向‘成本驱动’，商业化节奏取决于成本曲线下降速度。\par}
\paragraph{Citi} 阳光电源发布10kV商用固态变压器（SST），瞄准800V高压直流AI数据中心市场。峰值转换效率98.5\%，体积仅为传统变压器的三分之一，数据中心占地面积可减少70\%。五年研发投入年均超5亿元，300人团队，自研碳化硅高频拓扑。英伟达、亚马逊、微软代表到场，目标2026年底获美国UL认证。
{\small\color{refgray}数据：峰值转换效率98.5\%，比传统变压器高3-5个百分点\par}
{\small\color{refgray}数据：体积仅为传统油浸变压器三分之一，数据中心占地减少70\%\par}
{\small\color{refgray}数据：五年研发，年均资本开支超5亿元，300人团队\par}
{\small\color{refgray}数据：目标2026年底获美国UL认证\par}
{\small\color{refgray}边际变化：AI数据中心供电架构从‘多级变压整流’向‘一步到位直流转换’演进，固态变压器成为新品类。\par}
\paragraph{Citi} 软件支出环比改善，但增量资金高度集中在AI基础设施与数据层，传统应用软件仍在承压。2Q IT支出改善部分来自年初推迟项目交付，但改善分布极不均匀。现代数据架构是AI部署前提，Snowflake、MongoDB、Palantir等数据平台处于战略位置，传统应用软件面临更严格ROI审视。
{\small\color{refgray}数据：2Q IT支出环比改善，但增量资金集中于AI基础设施与数据层\par}
{\small\color{refgray}数据：现代数据架构是AI部署前提条件\par}
{\small\color{refgray}数据：传统应用软件面临整合不确定性与ROI审视\par}
{\small\color{refgray}数据：AI支出增长与软件全面复苏需明确切割\par}
{\small\color{refgray}边际变化：AI增量资金未溢出到传统软件品类，市场需区分‘受益于AI基础设施’与‘受益于软件支出回暖’。\par}
\paragraph{GS} 微软取消数据中心协议非需求见顶，电力才是真正约束。前微软AI转型总监指出，Agentic AI与企业内部未训练数据提供双重需求支撑，AI基础设施不存在过度建设风险。企业正将AI工作负载从云端拉回本地，原因在于成本、延迟与数据主权，DELL、HPE等企业计算厂商迎来结构性利好。
{\small\color{refgray}数据：Agentic AI推动工作负载从云端向本地迁移\par}
{\small\color{refgray}数据：企业采用混合计算策略：本地运行AI Agent，云端负责多代理编排\par}
{\small\color{refgray}数据：互联网数据仅占全球数据总量约三分之一\par}
{\small\color{refgray}数据：新一代CPU服务器（如Arm架构）单机柜性能提升，功耗降低约50\%\par}
{\small\color{refgray}边际变化：AI算力需求从‘集中训练’向‘混合推理’迁移，推理侧硬件需求可能被低估。\par}
\paragraph{GS} 本土半导体设备公司受益于本土存储与先进制程产能扩张，毛利率连续三个季度站上40\%，但新产品爬坡节奏慢于预期导致营收低于预测。CXMT与腾讯签署服务器DRAM长期供应协议，表明本土存储客户需求向高端应用迁移。2026年净利润预测下调5\%，但2027-2028年营收与毛利率预测上调。
{\small\color{refgray}数据：1Q26毛利率40.0\%，连续三个季度超过40\%\par}
{\small\color{refgray}数据：1Q26营收同比下降11\%、环比下降19\%，低于预期\par}
{\small\color{refgray}数据：CXMT与腾讯签署服务器DRAM长期供应协议\par}
{\small\color{refgray}数据：2026年净利润预测下调5\%，2027-2028年营收与毛利率预测上调\par}
{\small\color{refgray}边际变化：毛利率改善确定性高于营收恢复节奏，两者存在时间差；本土存储客户需求从消费级向服务器级升级。\par}
\paragraph{GS} 沪电股份泰国产能爬坡超预期，1Q26产能利用率超90\%，净销售额2.95亿元，超70\%数据通信海外客户完成工厂验证。二期扩产已启动，2Q26新产能投产。AI网络需求是产品结构升级的真正推手，公司在AI网络PCB市场占据领先地位，产品组合升级直接提升盈利能力。
{\small\color{refgray}数据：泰国工厂1Q26产能利用率超90\%\par}
{\small\color{refgray}数据：净销售额2.95亿元\par}
{\small\color{refgray}数据：超70\%数据通信海外客户完成工厂验证\par}
{\small\color{refgray}数据：二期扩产2Q26投产\par}
{\small\color{refgray}边际变化：海外产能从‘试水’进入‘满产’状态，关注焦点从‘能否建好’转向‘扩产速度能否跟上订单增长’。\par}
\paragraph{GS} 全球资金流向在7月第一周净流入约563亿美元，扭转前一周净流出，但驱动力量高度集中于科技板块（单周净流入约288.6亿美元）。韩国与台湾市场资金流入加速，GS策略团队认为‘记忆体超级周期’支撑盈利增长预期，但伴随高波动性。科技板块Z-score达4.50，属于统计极端值。
{\small\color{refgray}数据：全球股票基金单周净流入约563亿美元\par}
{\small\color{refgray}数据：科技板块单周净流入约288.6亿美元\par}
{\small\color{refgray}数据：科技板块四周累计流入约592亿美元\par}
{\small\color{refgray}数据：科技板块Z-score达4.50\par}
{\small\color{refgray}边际变化：资金流向从‘全面回暖’转向‘高度集中于科技’，记忆体超级周期成为核心驱动变量。\par}
\subsubsection*{关键数据}
\begin{itemize}
  \item 全球IT预算增速回升至+3.3\%，AI预算占比6.5\%
  \item 50\%的CIO认为AI支出挤占其他IT预算（3月40\%）
  \item 大族数控2Q26净利润同比增长294\%-362\%
  \item 东方电气签署10台50MW燃气轮机出口订单，毛利率超30\%
  \item 阳光电源SST转换效率98.5\%，体积减少70\%
  \item 沪电股份泰国产能利用率超90\%，二期扩产启动
  \item 联发科AI ASIC 2028年收入预计220亿美元
  \item 科技板块单周资金净流入约288.6亿美元，Z-score 4.50
\end{itemize}
\begin{figure}[htbp]
\centering
\includegraphics[width=0.88\linewidth]{figures/fig_027.jpg}
\caption{Figure 1 - IT Servies (Keane) - As we continue to monitor the evolution of IT budgets in the AI world within the IT services industry, a major trend we are seeing is that enterprise clients/CIOs are reallocating spend from the AI productivity gains to invest in AI-relate}
\end{figure}
\begin{figure}[htbp]
\centering
\includegraphics[width=0.88\linewidth]{figures/fig_032.jpg}
\caption{Figure 2 - Figure 2. HC: 1H26/2Q26 preliminary results snapshot © 2026 Citi Inc. No redistribution without Citi's written permission. Figure 3. HL: Revenue by segment and YoY growth, 1H26 © 2026 Citi Inc. No redistribution without Citi's written permission. Figure 4. P/E}
\end{figure}
\begin{figure}[htbp]
\centering
\includegraphics[width=0.88\linewidth]{figures/fig_035.jpg}
\caption{Figure 1 - Much interest from US potential buyers – According to Sungrow, representatives from NVIDIA, Amazon and Microsoft have also joined the event in Anhui, China today. We believe that these American companies are interested in the new product. This product seems li}
\end{figure}
\begin{figure}[htbp]
\centering
\includegraphics[width=0.88\linewidth]{figures/fig_044.jpg}
\caption{Exhibit 2 - Exhibit 2: Peers' correlation of 2027E P/E to 2027-28E avg. EPS growth Exhibit 3: E-Town Semis 12M forward P/E Exhibit 4: E-Town and peers' PEG}
\end{figure}
{\small\color{refgray}本节综合 12 篇研究；涉及 Bernstein、Citi、GS。}
\newpage
\subsection{AI基础设施需求的结构性迁移：从电信到数据中心}
\textbf{AI数据中心（AIDC）正取代传统电信成为光纤、CCL等基础设施需求的核心驱动力，推动相关材料与组件进入量价齐升的新周期。}
\subsubsection*{主流共识}
\begin{itemize}
  \item AI数据中心对光纤光缆的增量需求正在弥补电信市场的放缓，行业需求结构已发生实质性转移。
  \item CCL等电子材料受益于AI与高速运算需求扩散，涨价趋势从高端向大宗品类蔓延。
  \item CPO（共封装光学）技术商业化路径清晰，关键组件量产时间表明确，供应链加速成熟。
\end{itemize}
\subsubsection*{分歧与条件差异}
\begin{itemize}
  \item 中国移动光缆招标量同比下降6.6\%，但市场对价格走势存在分歧：部分观点认为折扣压力仍存，另一部分认为AIDC需求将推动报价中枢上移。
  \item CPO技术路线选择存在差异：光栅耦合（GC）与边缘耦合（EC）方案各有支持者，影响供应链格局。
  \item 小米AI能力（MiMo）的估值影响存在分歧：部分机构认为其日处理67万亿token的能力被低估，另一些则关注其汽车业务对估值的稀释。
\end{itemize}
\subsubsection*{机构观点}
\paragraph{GS} 小米AI能力被低估，MiMo日处理token量达67万亿，API请求量6月较1月增长27倍。SkyNomad增程SUV系列以可配置智能空间为卖点，纯电续航265-370公里领先同级，有望成为叙事与财务转折点。
{\small\color{refgray}数据：MiMo日处理token量67万亿\par}
{\small\color{refgray}数据：API请求量6月较1月增长27倍\par}
{\small\color{refgray}数据：SkyNomad纯电续航265-370公里\par}
{\small\color{refgray}数据：小米股价QTD跑赢恒生科技指数13个百分点\par}
{\small\color{refgray}边际变化：小米AI能力从被低估转向市场重新定价，SkyNomad产品逻辑从驾驶性能转向智能空间，拓展了估值叙事。\par}
\paragraph{JPM} 南亚塑胶电子材料板块营业利润率自2021年来首次突破20\%，CCL自3月中旬起提价15\%，AI需求扩散至大宗CCL品类。台塑集团二季度净利润640亿新台币创五年新高，电子材料是超预期核心。
{\small\color{refgray}数据：南亚塑胶电子材料营业利润率突破20\%\par}
{\small\color{refgray}数据：CCL价格自3月中旬起上调15\%\par}
{\small\color{refgray}数据：台塑集团Q2净利润640亿新台币\par}
{\small\color{refgray}数据：建滔集团年内已宣布6次涨价\par}
{\small\color{refgray}边际变化：电子材料利润率从个位数跳升至20\%以上，涨价范围从高端向大宗品类扩散，AI需求成为结构性驱动力。\par}
\paragraph{NOM} 中国移动光缆招标年化需求9230万芯公里，同比下降6.6\%，但光纤需求重心已从电信转向AI数据中心。现货价格明显上涨，预计招标报价将高于去年，行业定价权可能从买方转向卖方。
{\small\color{refgray}数据：年化需求9230万芯公里，同比-6.6\%\par}
{\small\color{refgray}数据：2025年招标加权均价53.85元/芯公里\par}
{\small\color{refgray}数据：光纤价格低于25元/芯公里\par}
{\small\color{refgray}数据：长飞光纤报价两周回调近40\%\par}
{\small\color{refgray}边际变化：光纤需求从电信驱动转向AI驱动，行业定价逻辑从量价齐跌转向量缩价升。\par}
\paragraph{NOM} 大立光CPO光纤阵列（FA）量产信心增强，已锁定一家核心客户，计划2026年9月试产、2027年中量产。公司主攻光栅耦合（GC）路线，不受康宁边缘耦合方案影响，同时储备模压玻璃棱镜方案。
{\small\color{refgray}数据：CPO FA量产时间表：2026年9月试产，2027年中量产\par}
{\small\color{refgray}数据：已锁定一家客户用于单排FA产品\par}
{\small\color{refgray}数据：目标价上调至6000新台币（基于2028年EPS 240新台币的25倍PE）\par}
{\small\color{refgray}数据：当前报价约3950新台币\par}
{\small\color{refgray}边际变化：CPO FA从概念验证进入客户锁定和试产排期阶段，商业化路径清晰，成为第二增长曲线的关键观测节点。\par}
\subsubsection*{关键数据}
\begin{itemize}
  \item MiMo日处理token量67万亿，API请求量6月较1月增长27倍
  \item 南亚塑胶电子材料营业利润率突破20\%，CCL提价15\%
  \item 中国移动光缆年化需求9230万芯公里，同比-6.6\%
  \item 大立光CPO FA量产时间表锁定2027年中
  \item 台塑集团Q2净利润640亿新台币创五年新高
  \item 长飞光纤报价两周回调近40\%
  \item 建滔集团年内已宣布6次涨价
  \item 小米股价QTD跑赢恒生科技指数13个百分点
\end{itemize}
\begin{figure}[htbp]
\centering
\includegraphics[width=0.88\linewidth]{figures/fig_086.jpg}
\caption{Exhibit 6 - Exhibit 8: Xiaomi MiMo API requests grew 27x in Jun vs. Jan, or 6-7x vs. Apr}
\end{figure}
\begin{figure}[htbp]
\centering
\includegraphics[width=0.88\linewidth]{figures/fig_088.jpg}
\caption{Figure 3 - Figure 5: Formosa group 2025 revenue mix}
\end{figure}
\begin{figure}[htbp]
\centering
\includegraphics[width=0.88\linewidth]{figures/fig_083.jpg}
\caption{Exhibit 1 - Exhibit 1: Comparable models to Xiaomi SkyNomad Exhibit 2: Exterior design of Xiaomi SkyNomad N90 Exhibit 3: Xiaomi SkyNomad N90 Max features a built-in pop-up rooftop tent Exhibit 4: SkyNomad offers configurable layouts which}
\end{figure}
\begin{figure}[htbp]
\centering
\includegraphics[width=0.88\linewidth]{figures/fig_089.jpg}
\caption{Figure 5 - Figure 5: Formosa group 2025 revenue mix Formosa Petrochemical (6505 TT) Formosa Chemicals \& Fiber (1326 TT) Nan Ya Plastics (1303 TT)}
\end{figure}
{\small\color{refgray}本节综合 4 篇研究；涉及 GS、JPM、NOM。}
\newpage
\subsection{消费与医疗：分化加剧，创新与定价权成核心变量}
\textbf{中国消费需求结构性改善尚未出现，运动品牌批发渠道持续承压；医疗领域基药目录扩容为创新药打开基层市场，但放量节奏取决于渠道执行；全球碳转型框架下，消费医疗行业碳评分领先，但转型计划可信度成为新焦点。}
\subsubsection*{主流共识}
\begin{itemize}
  \item 中国运动用品需求未见结构性改善，批发渠道数据持续疲弱，宝胜6月零售额同比下滑8.6\%，两年复合增速创两年新低。
  \item 全球碳定价方向明确，消费医疗行业在碳强度、减排雄心与执行能力上领先，但转型计划可信度成为资金配置的关键考量。
  \item Netflix已跨越用户增长驱动阶段，定价权与广告收入成为利润扩张新引擎，市场对参与度指标的过度担忧可能掩盖结构性变化。
\end{itemize}
\subsubsection*{分歧与条件差异}
\begin{itemize}
  \item 运动品牌需求：Bernstein认为批发渠道数据反映需求持续调整，但部分市场观点可能期待基数效应带来改善。
  \item Netflix估值：JPM认为当前报价过度反映短期担忧，下半年利润将加速；但市场对用户参与度下滑和提价退订风险仍有分歧。
  \item 基药目录影响：NOM强调创新药基层市场准入是真正增量，但部分观点可能认为数量扩容有限，放量效果待观察。
  \item 电池隔膜价格：UBS预期提价7-8\%将扭转行业趋势，但价格谈判落地时间和幅度仍存不确定性。
\end{itemize}
\subsubsection*{机构观点}
\paragraph{Bernstein} 中国运动用品需求未见结构性改善。宝胜6月零售额同比下滑8.6\%，两年复合增速-24.9\%，创两年新低，且连续第六年在该月份收入同比下降。裕元制造端6月收入同比-11.8\%，波动未缓解。批发渠道数据过滤直营与电商噪音，直接反映库存消化与终端动销，品牌方出货量难以真正回暖。
{\small\color{refgray}数据：宝胜6月零售额同比-8.6\%\par}
{\small\color{refgray}数据：两年复合增速-24.9\%\par}
{\small\color{refgray}数据：连续6年6月收入同比下降\par}
{\small\color{refgray}数据：裕元6月收入同比-11.8\%\par}
{\small\color{refgray}边际变化：6月负增长延续调整趋势，两年复合增速创两年新低，说明非基数问题，需求加速下滑。\par}
\paragraph{Bernstein} 全球碳转型框架下，消费医疗行业碳评分领先，但转型计划可信度成为关键。Bernstein与SG联合构建四维评分体系（碳强度、趋势、雄心、执行），对STOXX 600排名。SFDR 2.0将新增“转型”类别，允许高排放行业通过可信计划获得资本。未管理的碳排放正变成隐性负债。
{\small\color{refgray}数据：2024年全球平均气温超工业化前1.55°C\par}
{\small\color{refgray}数据：剩余碳预算仅够约3年\par}
{\small\color{refgray}数据：SFDR 2.0新增转型类别\par}
{\small\color{refgray}边际变化：碳从ESG披露指标变为财务变量，市场焦点从历史排放转向转型计划可信度。\par}
\paragraph{JPM} Netflix已跨越用户增长驱动阶段，定价权与广告收入推动利润扩张。全球DAU同比-6\%，但内部质量参与度指标创历史新高。美国提价预计年化增量17亿美元，广告收入2026年预计31亿美元，2027年达53亿美元。下半年费用增速从16\%下降，利润将明显加速。
{\small\color{refgray}数据：全球DAU同比-6\%\par}
{\small\color{refgray}数据：美国提价年化增量17亿美元\par}
{\small\color{refgray}数据：2026年广告收入预计31亿美元\par}
{\small\color{refgray}数据：2027年广告收入预计53亿美元\par}
{\small\color{refgray}边际变化：参与度与收入增长相关性减弱，市场过度关注DAU下滑，忽视定价权与广告业务的结构性贡献。\par}
\paragraph{NOM} 新版基药目录（2026版）总品种从685扩至794，净增109种，真正增量在创新药。GLP-1、肿瘤单抗、靶向药首次大规模纳入，糖尿病和肿瘤最受益。石药丁苯酞、恒瑞碘佛醇、信达贝伐珠单抗和阿达木单抗、荣昌泰它西普等获基层市场准入。目录9月1日实施，放量逻辑不同于医保谈判。
{\small\color{refgray}数据：总品种从685扩至794\par}
{\small\color{refgray}数据：净增109种\par}
{\small\color{refgray}数据：GLP-1、肿瘤单抗首次纳入\par}
{\small\color{refgray}数据：信达贝伐珠单抗和阿达木单抗同时纳入\par}
{\small\color{refgray}边际变化：基药目录时隔8年更新，创新药大规模进入，基层用药结构从“保基本”向“有选择”切换。\par}
\paragraph{UBS} 恩捷股份二季度业绩超预期，量价修复同步。隔膜出货41亿平方米，同比+58\%，环比+17\%，超指引。单位净利润从0.13元/平米回升至0.15元/平米。行业价格谈判接近完成，预期提价7-8\%（约0.05-0.06元/平米），若落地将扭转2024年以来价格下行趋势。
{\small\color{refgray}数据：二季度归母净利润中值6.1亿元\par}
{\small\color{refgray}数据：环比+115\%\par}
{\small\color{refgray}数据：隔膜出货41亿平米，同比+58\%\par}
{\small\color{refgray}数据：单位净利润0.15元/平米\par}
{\small\color{refgray}边际变化：出货量超预期且价格谈判接近完成，市场此前对隔膜环节定价过于保守，提价预期将改善下半年盈利弹性。\par}
\subsubsection*{关键数据}
\begin{itemize}
  \item 宝胜6月零售额同比-8.6\%，两年复合增速-24.9\%
  \item Netflix全球DAU同比-6\%，内部质量参与度指标创历史新高
  \item 新版基药目录净增109种，GLP-1、肿瘤单抗首次纳入
  \item 恩捷二季度隔膜出货41亿平米，同比+58\%，超指引
  \item 隔膜行业预期提价7-8\%
  \item 2024年全球平均气温超工业化前1.55°C
  \item 剩余碳预算仅够约3年
\end{itemize}
\begin{figure}[htbp]
\centering
\includegraphics[width=0.88\linewidth]{figures/fig_004.jpg}
\caption{EXHIBIT 1 - EXHIBIT 3: Last month's revenue was below Pre-COVID levels Pou Sheng Monthly Revenue}
\end{figure}
\begin{figure}[htbp]
\centering
\includegraphics[width=0.88\linewidth]{figures/fig_074.jpg}
\caption{Figure 3 - Figure 3: NFLX DAU Growth – GLOBAL – 1/1/23 to 6/30/26 Figure 4: NFLX DAU Growth – UCAN – 1/1/23 to 6/30/26 \#\# DAU Takeaways:}
\end{figure}
\begin{figure}[htbp]
\centering
\includegraphics[width=0.88\linewidth]{figures/fig_001.jpg}
\caption{EXHIBIT 1 - EXHIBIT 1: Ambition ratings 1 to 5 (the best) EXHIBIT 2: Execution ratings 1 to 5 (the best) \#\# Company-level carbon signals: ambition \& execution ratings \#\# The scope}
\end{figure}
\begin{figure}[htbp]
\centering
\includegraphics[width=0.88\linewidth]{figures/fig_005.jpg}
\caption{EXHIBIT 1 - EXHIBIT 3: Last month's revenue was below Pre-COVID levels Pou Sheng Monthly Revenue EXHIBIT 4: The June 2026 result makes it 6 years of consecutive revenue declines in the month Pou Sheng Monthly Revenue - June}
\end{figure}
{\small\color{refgray}本节综合 5 篇研究；涉及 Bernstein、JPM、NOM、UBS。}
\newpage
\subsection{中国汽车与工业：出口驱动结构升级，内需筑底与格局重塑}
\textbf{中国汽车行业呈现内需疲弱与出口强劲的显著分化，新能源渗透率突破62\%但龙头份额承压，行业正从价格竞争转向技术与品牌竞争，出口结构从量增转向质变，欧洲成为关键验证市场。}
\subsubsection*{主流共识}
\begin{itemize}
  \item 中国新能源乘用车渗透率6月突破62\%，纯电增速重新超过插混，但国内零售连续三个月两位数下滑，内需处于筑底阶段。
  \item 出口成为行业核心增长支柱，上半年乘用车出口同比增长72\%，EV出口增速是燃油车的四倍，欧洲已取代亚洲成为中国品牌第一大出口目的地。
  \item 国内价格战趋缓，平均折扣从3月17.5\%收窄至6月底15.7\%，但电池等BOM成本上升压缩车企降价空间，竞争焦点转向技术和品牌。
\end{itemize}
\subsubsection*{分歧与条件差异}
\begin{itemize}
  \item 市场增量流向分歧：Citi数据显示比亚迪份额在渗透率上升期反而下滑，增量被零跑、吉利等品牌瓜分；而GS认为本土科技与豪华品牌（小鹏、极氪）正在填补价格带空白。
  \item 欧洲豪华品牌在华表现分化：奔驰、宝马因产品空窗期份额损失较大，而奥迪、沃尔沃凭借本土化新品（AUDI E7X、XC70 PHEV）表现更具韧性。
  \item 霍尼韦尔分拆后利润率修复路径存在不确定性：Bernstein预测EBIT利润率从2025年12.4\%跃升至2027年20.1\%，但实际执行节奏可能受分拆一次性费用和需求端影响。
\end{itemize}
\subsubsection*{机构观点}
\paragraph{Bernstein} 霍尼韦尔分拆后RemainCo估值并不激进，20倍远期市盈率对应每股约204美元，叠加Quantinuum持股每股约39美元，目标价243美元。利润率修复是核心看点：EBIT利润率从2025年12.4\%爬升至2027年20.1\%，但需区分分拆一次性费用与真实经营改善。上行风险来自楼宇自动化和Forge平台获客，下行风险包括能源价格稳定、LNG项目延迟及分拆组织阵痛。
{\small\color{refgray}数据：目标价243美元，较当前报价上行约7-9\%\par}
{\small\color{refgray}数据：2026年预计总收入约204亿美元，2027年约211亿美元\par}
{\small\color{refgray}数据：EBIT利润率：2025年12.4\% → 2026年17.7\% → 2027年20.1\%\par}
{\small\color{refgray}数据：Quantinuum持股估值约39美元/股\par}
{\small\color{refgray}边际变化：分拆后市场关注点从‘瘦身后还剩什么’转向Quantinuum股权价值与利润率修复路径，Bernstein首次将量子计算资产单独估值。\par}
\paragraph{Citi} 6月中国新能源乘用车批发148万辆，渗透率62.8\%创新高，但前五大厂商CR5仅57.1\%，同比降1.4个百分点。比亚迪份额26.8\%，同比降3.6个百分点，环比增速5\%跑输行业10\%均值，增量被零跑（同比+95\%）、吉利等品牌瓜分。插混增速放缓至+6\%，纯电重新成为增长主力（同比+27\%）。市场结构从龙头集中转向多品牌分化。
{\small\color{refgray}数据：6月新能源批发148万辆，同比+19\%，环比+10\%\par}
{\small\color{refgray}数据：渗透率62.8\%，同比+13个百分点\par}
{\small\color{refgray}数据：比亚迪市占率26.8\%，同比-3.6个百分点\par}
{\small\color{refgray}数据：零跑销量9.3万辆，同比+95\%，市占率6.3\%（+2.4个百分点）\par}
{\small\color{refgray}边际变化：渗透率突破60\%后，市场增量并未流向龙头，而是被更多品牌瓜分，竞争格局从‘一超多强’转向‘多极分化’。\par}
\paragraph{GS} 中远海能二季度净利润约23亿元，环比+7\%，远低于招商轮船的+50\%，主因7艘油轮（3艘VLCC、3艘LR2、1艘Panamax）自2月28日霍尔木兹海峡关闭后被困波斯湾，合计414天无法收取滞期费。6月海峡临时开放后船只已全部驶出，预计三季度起业绩追赶同业。VLCC二季度平均TCE约12.5万美元/天（环比+31\%），全年均价或维持15万美元/天附近。
{\small\color{refgray}数据：二季度净利润约23亿元，环比+7\%\par}
{\small\color{refgray}数据：7艘油轮合计被困414天，最长88天\par}
{\small\color{refgray}数据：VLCC二季度平均TCE约12.5万美元/天（环比+31\%）\par}
{\small\color{refgray}数据：2026年预测PE约4.5倍，股息率预计11.3\%\par}
{\small\color{refgray}边际变化：地缘事件导致的一次性收入损失正在消除，市场关注点从‘利润差距’转向‘三季度追赶同业’及运费率走向。\par}
\paragraph{GS} 6月中国乘用车零售160万辆，同比-23.2\%，新能源渗透率62.9\%。德国合资品牌承受最大份额损失：奔驰零售2.48万辆（同比-41.8\%），份额1.5\%（-49bps）；大众零售10.35万辆（同比-37.2\%），份额6.5\%（-156bps）。两者燃油车占比分别高达99.9\%和97.7\%，新产品（探岳L PHEV）仍处爬坡期。奥迪受益于AUDI E7X上市，份额损失较小（-39bps）。
{\small\color{refgray}数据：6月零售160万辆，同比-23.2\%\par}
{\small\color{refgray}数据：新能源渗透率62.9\%\par}
{\small\color{refgray}数据：奔驰零售2.48万辆，同比-41.8\%，份额1.5\%\par}
{\small\color{refgray}数据：大众零售10.35万辆，同比-37.2\%，份额6.5\%\par}
{\small\color{refgray}边际变化：燃油车零售同比-39\%加速萎缩，德国合资品牌产品切换期‘旧线加速调整、新品未成规模’的时间差成为核心矛盾。\par}
\paragraph{JPM} 6月国内乘用车销量同比-26\%至150万辆，出口同比+80\%至90万辆，上半年出口超440万辆（同比+72\%）。欧洲已占中国品牌出口37\%，超越亚洲成为第一大目的地。中国品牌在欧洲NEV注册份额从2025年约11\%升至上半年约17\%，5月达19\%。驱动因素包括产品竞争力、多动力总成选择、更快投放节奏及品牌接受度提升。国内价格战趋缓，平均折扣从3月17.5\%收窄至6月底15.7\%。
{\small\color{refgray}数据：6月国内乘用车销量150万辆，同比-26\%\par}
{\small\color{refgray}数据：6月出口90万辆，同比+80\%\par}
{\small\color{refgray}数据：上半年出口超440万辆，同比+72\%\par}
{\small\color{refgray}数据：欧洲占中国品牌出口37\%，NEV注册份额17\%\par}
{\small\color{refgray}边际变化：出口从‘量增’转向‘质变’，欧洲成为关键验证市场；国内需求不会V型反转，但价格战趋缓为长期健康市场铺路。\par}
\paragraph{NOM} 上半年国内乘用车零售870万辆，同比-20.1\%，创2020年以来最低。6月零售160万辆，同比-23.2\%，连续三个月两位数下滑。出口成为唯一增长支柱：6月出口90.5万辆（同比+80.3\%），其中EV出口50万辆（同比+154\%），增速是燃油车四倍。政府‘反内卷’政策叠加BOM成本上升，价格战空间收窄，行业从‘卷价格’转向‘卷技术’。下半年国内环比改善可期，但全年仍将录得两位数下滑。
{\small\color{refgray}数据：上半年国内零售870万辆，同比-20.1\%\par}
{\small\color{refgray}数据：6月零售160万辆，同比-23.2\%\par}
{\small\color{refgray}数据：6月出口90.5万辆，同比+80.3\%\par}
{\small\color{refgray}数据：EV出口50万辆，同比+154\%\par}
{\small\color{refgray}边际变化：国内零售连续三个月两位数下滑，出口成为行业唯一增长支柱，EV出口增速是燃油车四倍，行业转型从‘价格竞争’转向‘技术品牌竞争’。\par}
\subsubsection*{关键数据}
\begin{itemize}
  \item 6月中国新能源乘用车渗透率62.8\%，同比+13个百分点
  \item 比亚迪6月市占率26.8\%，同比-3.6个百分点
  \item 上半年中国乘用车出口超440万辆，同比+72\%
  \item 6月EV出口50万辆，同比+154\%，增速是燃油车四倍
  \item 欧洲占中国品牌出口37\%，NEV注册份额17\%
  \item 国内平均折扣从3月17.5\%收窄至6月底15.7\%
  \item 霍尼韦尔分拆后EBIT利润率目标：2025年12.4\%→2027年20.1\%
  \item 中远海能7艘油轮合计被困414天，二季度净利润环比+7\%低于同业+50\%
\end{itemize}
\begin{figure}[htbp]
\centering
\includegraphics[width=0.88\linewidth]{figures/fig_011.jpg}
\caption{Figure 1 - Figure 1. NEV Wholesale Volumes (Units) of Major OEMs © 2026 Citi Inc. No redistribution without Citi's written permission. Figure 2. BEV/PHEV Sector Wholesales © 2026 Citi Inc. No redistribution without Citi's written permission. \#\# Appendix A-1}
\end{figure}
\begin{figure}[htbp]
\centering
\includegraphics[width=0.88\linewidth]{figures/fig_046.jpg}
\caption{Exhibit 1 - Exhibit 1: China domestic PV retail volume dashboard As of Jun-26 \#\# Historicals vs. recent R12M retail volume}
\end{figure}
\begin{figure}[htbp]
\centering
\includegraphics[width=0.88\linewidth]{figures/fig_051.jpg}
\caption{Exhibit 1 - Exhibit 1: European premium is defending its share better than mass peers, but the loss is widening. VW Group mass brands remain the most under pressure among international brands. Market share change and absolute share in China, a}
\end{figure}
\begin{figure}[htbp]
\centering
\includegraphics[width=0.88\linewidth]{figures/fig_078.jpg}
\caption{Figure 1 - Figure 1: Chinese brands NEV market share in Europe (annual) Figure 2: Chinese brands NEV market share in Europe (monthly) Figure 3: Chinese brands NEV market share in WE5 (annual)}
\end{figure}
{\small\color{refgray}本节综合 7 篇研究；涉及 Bernstein、Citi、GS、JPM、NOM。}
\newpage
\subsection{航天与新能源：成本拐点与汇率杠杆}
\textbf{中国商业航天在长征十号B网捕回收成功后进入成本拐点，而三星生物制剂二季度盈利超预期主要来自韩元汇率顺风，而非运营改善。}
\subsubsection*{主流共识}
\begin{itemize}
  \item 火箭一级回收复用是降低发射成本的核心路径，占总成本60-65\%。
  \item 三星生物制剂产能利用率维持高位，1-4号工厂满产，5号工厂爬坡。
  \item 中国商业航天发射成本目前约6万元/公斤，远高于猎鹰9号。
\end{itemize}
\subsubsection*{分歧与条件差异}
\begin{itemize}
  \item 网捕路线是否优于垂直着陆：NOM认为网捕去掉着陆腿可提升载荷10-15\%，但SpaceX已验证高频复用可靠性。
  \item 三星生物制剂盈利驱动因素：市场关注罢工影响，但GS认为汇率才是二季度超预期主因。
\end{itemize}
\subsubsection*{机构观点}
\paragraph{GS} 三星生物制剂二季度CDMO收入同比+30\%至1.314万亿韩元，调整后净利润同比+41\%至4700亿韩元，高于市场预期。核心驱动是韩元汇率：二季度美元/韩元平均1500，较公司假设的1400高出7\%，汇率每变动2.5\%营业利润同向变动5\%，仅汇率贡献约14\%利润增量。5月罢工对二季度无实质影响，潜在扰动推迟至三季度且部分可恢复。6月新增订单（亚洲6290万、美国6900万、欧洲3430万美元）强化中期盈利信心。
{\small\color{refgray}数据：CDMO收入同比+30\%至1.314万亿韩元\par}
{\small\color{refgray}数据：调整后净利润同比+41\%至4700亿韩元\par}
{\small\color{refgray}数据：美元/韩元平均1500 vs 公司假设1400\par}
{\small\color{refgray}数据：汇率每变动2.5\%营业利润同向变动5\%\par}
{\small\color{refgray}边际变化：市场此前聚焦罢工风险，GS指出汇率才是二季度超预期主因，且罢工影响被高估。\par}
\paragraph{NOM} 长征十号B于2026年7月10日实现全球首次网捕式火箭一级回收，中国成为继美国后第二个掌握该技术的国家。网捕路线将着陆腿的几吨死重转移到回收船，使有效载荷提升10-15\%，且海上甲板晃动对网捕影响小。当前中国发射成本约6万元/公斤，长十B目标降至2万元/公斤（约2800美元/公斤）。与SpaceX猎鹰9号（载荷比4.2\%、最高复用34次）相比，差距不在单次回收，而在高频可靠复用。
{\small\color{refgray}数据：火箭一级占总成本60-65\%\par}
{\small\color{refgray}数据：网捕去掉着陆腿使有效载荷提升10-15\%\par}
{\small\color{refgray}数据：当前中国发射成本约6万元/公斤\par}
{\small\color{refgray}数据：长十B目标成本2万元/公斤（约2800美元/公斤）\par}
{\small\color{refgray}边际变化：从技术演示转向成本拐点，中国首次实现轨道级火箭回收，路径选择不同于SpaceX。\par}
\subsubsection*{关键数据}
\begin{itemize}
  \item 三星生物制剂二季度CDMO收入同比+30\%至1.314万亿韩元
  \item 汇率每变动2.5\%营业利润同向变动5\%
  \item 长征十号B实现全球首次网捕式火箭一级回收
  \item 中国当前发射成本约6万元/公斤，目标降至2万元/公斤
  \item 猎鹰9号助推器最高复用34次，载荷比4.2\%
\end{itemize}
\begin{figure}[htbp]
\centering
\includegraphics[width=0.88\linewidth]{figures/fig_070.jpg}
\caption{Exhibit 1 - Exhibit 1: Avg USD/KRW at 1,500 in 2Q26 versus 1,400 factored in company guidance Exhibit 2: Samsung Biologics is trading at a 31x 12m fwd P/E Data as of Jul 6, 2026 Exhibit 3: ...versus Lonza at 30x 12m fwd P/E Data as of Jul}
\end{figure}
\begin{figure}[htbp]
\centering
\includegraphics[width=0.88\linewidth]{figures/fig_071.jpg}
\caption{Exhibit 1 - Exhibit 4: We see an upward guidance revision for Samsung Bio more likely in 3Q}
\end{figure}
\begin{figure}[htbp]
\centering
\includegraphics[width=0.88\linewidth]{figures/fig_072.jpg}
\caption{Exhibit 1 - Exhibit 4: We see an upward guidance revision for Samsung Bio more likely in 3Q}
\end{figure}
{\small\color{refgray}本节综合 2 篇研究；涉及 GS、NOM。}
\newpage
\section{辅助信号｜战略咨询}
本板块摘要：麦肯锡两份报告分别聚焦全球宏观环境的结构性张力与6G/AI产业的技术路径分歧。宏观方面，IMF上调增长预期但贸易摩擦滞后效应未消，美国消费与就业数据分化、中国消费结构变化值得关注。产业方面，AI在MWC上无处不在但可持续价值路径不明，6G定义尚未对齐，运营商在AI-RAN和GPUaaS上寻求新利润点。
\subsection{全球宏观的结构性张力：增长韧性掩盖贸易摩擦滞后效应}
\textbf{IMF上调全球增长预期至3.0\%，但麦肯锡指出这一上调主要受关税前抢跑、美元走弱和财政扩张的暂时支撑，而非贸易摩擦已消化。美国消费撑起3\%GDP但就业数据大幅下修（5月从+14.4万修正为+1.9万，6月从+14.7万修正为+1.4万，合计下修25.8万人），通胀黏性未消（6月CPI同比2.7\%，核心2.9\%）。中国Q2增长5.2\%但消费与贸易结构悄然变化。关税收入可作为衡量真实冲击的替代指标。}
\begin{itemize}
  \item IMF将2025年全球增长预期上调至3.0\%，但强调比4月2日前基线低约0.2个百分点，上调主因关税前抢跑、美元走弱和财政扩张，而非贸易摩擦本身已消化。
  \item 美国Q2 GDP环比年化增长3\%，主要靠家庭消费和净出口改善，但5月和6月非农就业合计下修25.8万人，就业市场实际强度远低于初值。
  \item 美国6月CPI同比升至2.7\%，核心CPI年化升至2.9\%，通胀不确定性未消退，美联储以分歧投票维持利率不变。
  \item 中国Q2增长5.2\%，消费与贸易结构出现边际变化，需关注内需对增长的支撑力度。
  \item 关税收入可作为衡量贸易摩擦真实冲击的替代指标：当关税收入持续增长、核心通胀开始反映关税传导、制造业PMI持续低于50时，滞后效应才进入显性阶段。
\end{itemize}
\subsection{6G与AI产业路径分歧：AI无处不在但可持续价值路径不明，6G定义尚未对齐}
\textbf{MWC 2025上AI渗透到每一场发布，但行业尚未回答什么能留下、什么能产生可持续价值。运营商在消费者端推出AI Phone、个人AI助手，在企业端推出AI数字助手和GPU即服务，试图从AI-RAN和GPUaaS中获取新利润。6G讨论聚焦AI原生和感知，但各方对定义尚未对齐。卫星直连手机从概念走向现实，全球91家运营商已签约。}
\begin{itemize}
  \item AI和生成式AI在MWC上无处不在，但哪些应用会真正留下来、哪些只是展台概念，行业尚无共识，运营商既兴奋又焦虑。
  \item 运营商在消费者端推出AI Phone（如德意志电信Magenta AI Phone搭载Perplexity）、个人AI助手（如SKT Astar预计2025年北美上线），在企业端推出AI数字助手和GPU即服务，试图利用靠近用户的低延迟优势。
  \item 6G讨论聚焦AI原生和感知，但各方对定义尚未对齐，行业仍在探索技术路径和标准化方向。
  \item 卫星直连手机从概念走向现实，全球91家运营商与卫星提供商签约，覆盖超60\%移动用户，欧盟投入100亿欧元打造自主卫星星座。
  \item 运营商在AI-RAN和GPUaaS上的布局，本质上是试图从过去十年未能从视频流量暴增和5G部署中分到足够利润的困境中寻找新出路。
\end{itemize}
\newpage
\section{辅助信号｜智库/国际机构}
本板块聚焦亚洲开发银行、IMF、经合组织、世界贸易组织等机构的最新研究，提炼出三个核心主题：全球供应链的结构性重组与竞争格局、新兴市场与前沿经济体的治理与效率瓶颈、以及绿色转型与可持续金融的落地挑战。这些报告提供了超越市场共识的实证视角，对宏观判断、产业战略和资产叙事具有辅助信号价值。
\subsection{全球供应链的结构性重组与竞争格局}
\textbf{全球供应链并未系统性脱钩，而是在经历产业、国家和企业所有权层面的静默重组；同时，AI和清洁技术等前沿领域的竞争格局正从下游向上游转移，结构性风险加剧。}
\begin{itemize}
  \item 经合组织最新报告显示，以实际值衡量，2024年全球价值链贸易占GDP比重仍维持在17\%的历史高位附近，并未出现系统性下降。名义值与实际值的背离源于2021-2022年能源价格暴涨，而非生产网络扩张。所谓“脱钩”更多是局部和战略性的，而非系统性现象。
  \item 经合组织报告指出，AI模型开发层竞争加剧（语言模型开发者从9家增至47家，活跃模型从22个增至453个），但上游芯片制造、云计算基础设施和数据资源已高度集中。模型层的竞争态势能否持续，取决于上游关键投入品的可及性，结构性垄断风险正在加剧。
  \item 布鲁盖尔研究所认为，欧盟《产业加速法案》设定的2035年制造业占GDP 20\%的硬性目标缺乏宏观环境逻辑，可能扭曲资源配置。其“欧洲制造”产地规则将推高成本、延迟脱碳，并面临WTO挑战。报告建议放弃该目标，转而关注脱碳进度和产业竞争力指标。
  \item IMF工作论文揭示，中东和中亚地区的生产率增长瓶颈不在资源禀赋，而在竞争不足。关税越高，企业定价能力增长越快，生产率差距越大。制造业的markup正在快速追赶服务业，竞争不足正从结构性问题演变为系统性特征。
  \item 世界贸易组织部长级会议将化石燃料补贴改革从联合声明推进到产出文件阶段，建立了透明化、扰动应对和识别有害补贴三大支柱的实操工具箱。透明机制进入贸易政策审议的“提问清单”阶段，为跨国企业评估能源成本和竞争环境提供了新的公开信息来源。
\end{itemize}
\begin{figure}[htbp]
\centering
\includegraphics[width=0.88\linewidth]{figures/fig_157.jpg}
\caption{Figure 2 - Before the global financial crisis, both nominal and real measures show a strong increase in trade via GVCs, confirming their rapid expansion during the period of “hyperglobalisation” (Subramanian and Kessler, 2014[4]). After 2011, however, the two series dive}
\end{figure}
\begin{figure}[htbp]
\centering
\includegraphics[width=0.88\linewidth]{figures/fig_143.jpg}
\caption{Figure 1 - Empirical evidence suggests continuous innovation and declining prices, especially in the AI model development segment. Despite significant entry costs and the presence of a few tech giants in the AI model development segment \$\textasciicircum{}\{1\}\$ , the performance of Genera}
\end{figure}
\subsection{新兴市场与前沿经济体的治理与效率瓶颈}
\textbf{新兴市场和发展中经济体的核心瓶颈往往不是资源或资金不足，而是治理碎片化、执行赤字和资源配置效率问题，这决定了政策效果和投资环境的真实边界。}
\begin{itemize}
  \item IMF对纳米比亚社会支出的分析显示，该国教育、医疗和社会救助支出水平并不低，但产出效果滞后。教育支出中大部分锁定在教职工工资等经常性支出，而资本性支出比例偏低，且资源向高等教育倾斜，基础教育生均支出低于新兴市场平均水平。这是典型的“效率缺口”案例。
  \item 经合组织为马耳他制定的技能战略报告指出，该国技能体系的核心瓶颈不是资金或培训机会，而是治理碎片化。决策权分散在多个部委和机构之间，没有一个实体拥有跨部门协调的正式授权，导致数据不共享、需求预测与培训供给脱节。
  \item 经合组织对土耳其反海外贿赂体系的跟进报告显示，在71项建议中仅完全落实10项，49项未采取任何实质性行动。三大优先建议（吹哨人保护立法、企业刑事责任条款、海外贿赂指控传递机制）全部落空，形成“制度性执行赤字”。
  \item IMF技术援助报告揭示，伊拉克央行已成功开发宏观环境基础工具（MFT），但该工具尚未进入央行管理层决策会议，也未能与财政部中期预算信息有效对接。一个技术上成熟的预测工具仍停留在“实验室”阶段，数据质量与跨部门协调是模型落地的真正瓶颈。
  \item IMF评估指出，多米尼加共和国的财政透明度短板不在报表格式，而在边界扩展不足。公共信托产品、长期购电协议、养老金缺口和国企财务状况等关键领域数据缺失或分析流于形式，大量政府控制的资金和风险游离在正式监督之外。
\end{itemize}
\begin{figure}[htbp]
\centering
\includegraphics[width=0.88\linewidth]{figures/fig_105.jpg}
\caption{Figure 1 - 4. The paper is organized as follows. Section B provides the conceptual framework and situates social spending within the broader fiscal context. Section C assesses the education sector. Section D examines the health sector. Section E analyzes social protectio}
\end{figure}
\begin{figure}[htbp]
\centering
\includegraphics[width=0.88\linewidth]{figures/fig_138.jpg}
\caption{Figure 1 - Such a strategy is essential, given that skills policies are located at the intersection of education, labour market, industrial and other policy domains (OECD, 2019[12]). A coherent skills strategy can provide a clear roadmap for achieving Malta Vision 2050, }
\end{figure}
\subsection{绿色转型与可持续金融的落地挑战}
\textbf{绿色和可持续金融的规模扩张背后，存在标准不统一、流动性不足和激励机制错位等结构性障碍，这些因素决定了资本能否有效流向真正需要长期资金的项目和地区。}
\begin{itemize}
  \item 经合组织与卢森堡交易所的联合调研显示，全球绿色、社会和可持续债券中只有17\%来自发展中国家（剔除中国后降至5\%）。投资者兴趣真实，但实际配置仍集中在发达市场，流动性是发展中国家GSSS债券进入主流机构可投资清单的最大隐性门槛。
  \item 经合组织水安全融资报告指出，2015-2024年间全球贴标“水”或“蓝色”的债券累计发行约1.6万亿美元，但标签缺乏统一标准，单只债券平均标签数达1.27个。发行主体高度集中于主权和超主权机构，真正需要长期资金的地方水务公司和项目公司发行量有限。
  \item 亚洲开发银行简报强调，Bond Connect已从市场准入选项演变为覆盖交易、托管、外汇、衍生品和回购的完整跨境生态。到2025年已允许境外投资者参与买断式回购业务，并通过Swap Connect提供利率互换工具。决定资金是否长期停留的关键是基础设施兼容性，而非额度限制。
  \item IMF工作论文指出，稳定币在新兴市场的角色取决于汇率失衡程度。当失衡较低时，稳定币改善外汇准入和价格发现；当失衡较高时，其公共信号可能成为协调资本外逃的放大器。玻利维亚2024年允许虚拟资产交易后，USDT/BOB价格迅速成为平行美元的实际参考指标。
  \item 亚洲开发银行对兴都库什-喜马拉雅地区住房的研究显示，2000-2015年间地震高风险区域人口增长21\%，比低风险区域快3.4个百分点。抗震成本并非不可承受（仅增加5\%-10\%），但政策、资金与实践三者脱节，导致最脆弱群体既住不起安全房，也住不上抗震房。
\end{itemize}
\begin{figure}[htbp]
\centering
\includegraphics[width=0.88\linewidth]{figures/fig_152.jpg}
\caption{Figure 2 - \#\# FINDING 1: When deciding to invest in public sector GSS bonds (over conventional bonds with a similar risk profile), investors are driven by these instruments' alignment with institutional investment strategy and sustainability goals as well as their potent}
\end{figure}
\begin{figure}[htbp]
\centering
\includegraphics[width=0.88\linewidth]{figures/fig_147.jpg}
\caption{Figure 2 - Unlike GSS bonds, which allocate proceeds to specific projects, “sustainability-linked” bonds (SLBs) are general-purpose bonds, whose financial terms are tied to the achievement of environmental and social performance targets (OECD, 2024[9]). In other words, t}
\end{figure}
\subsection{宏观政策与金融稳定}
\textbf{用于校准利率、通胀、财政约束和金融稳定叙事。}
\begin{itemize}
  \item 国际研究争端解决正在经历一场静默的转向。仲裁成本飙升、周期漫长，调解作为一种低成本、高效率的替代方案，理论上备受推崇。但亚洲开发银行研究所最新政策简报指出，调解在读者与国家间的争端中远未普及，核心障碍并非技术或法律空白，而是政府自身的运作逻辑。
  \item 一份2025年5月完成、2026年1月发布的IMF技术援助报告，揭示了一个多数市场观察者容易忽略的事实：在非洲南部小国史瓦帝尼，非银行机构的资产规模已占据整个资金体系的66.2\%，但此前长期游离在官方货币与资金统计框架之外。报告的核心判断是，IMF正在帮助该国央行将养老金、保险、...
  \item 伊拉克宏观环境长期被一个问题困扰：高度依赖石油收入，油价波动直接冲击财政、货币和国际收支。中央银行需要一套能预测这些冲击、支持政策决策的分析框架。IMF的技术援助报告给出的结论很直接——工具已经有了，真正的挑战不在模型本身，而在团队能否持续用它、数据能否支撑它、央行能否与财政部真...
  \item 在固定汇率制度下，货币政策独立性天然受到制约。但这份IMF最新工作论文揭示了一个更具体的结论：纳米比亚的银行存贷款利率变动，主导力量来自南非储备银行的政策利率，而非本国央行的回购利率。这不仅是学术验证，更对任何实行货币局制度或紧密盯住汇率的地区，提供了一个关于政策传导边界的现实参...
  \item 一份来自国际货币产品组织（IMF）的最新工作论文，系统拆解了美国回购市场利差的驱动因素。结论并不复杂，但含义值得细品：准备金依然是回购市场最核心的稳定器，但它的效果并非在所有市场环境下都相同。真正决定回购利率走向的，是准备金、交易商资产负债表使用率、对冲产品杠杆与国债发行这四个变...
  \item 2025年3月，联合国统计委员会与国际货币产品组织同步发布了国民账户体系和国际收支手册的最新版本。对大多数国家的统计机构而言，这不仅仅是一次技术升级，而是一场需要跨部门协调、用户参与和资源重新配置的系统工程。IMF统计部门最新发布的《SNA/BPM实施规划手册》给出了一个清晰的信...
  \item 理解资金体系之外的宏观环境活动，对政策制定者而言从来不是可选项。但加密货币挖矿——这个不依赖传统资金中介、直接创造新资产的入口——其真实规模在全球范围内几乎是一团迷雾。IMF最新工作论文提出了一种全新的测量思路：通过追踪全球主要产地的矿机出口数据，来反推各国的挖矿活动分布。这个方...
\end{itemize}
\begin{figure}[htbp]
\centering
\includegraphics[width=0.88\linewidth]{figures/fig_120.jpg}
\caption{Figure 1 - \#\# 5. The close tracking of the BoN's policy rate with that of the SARB has changed in recent years. Prior to 2021, the BoN repo rate closely tracked the SARB policy rate and was typically set at a modest premium, reflecting the BoN's assessment of domestic \#\#}
\end{figure}
\begin{figure}[htbp]
\centering
\includegraphics[width=0.88\linewidth]{figures/fig_125.jpg}
\caption{Figure 1 - Figure 1: Spreads Are Largely Driven by Reserve Dynamics and Dealer Balance Sheet Usage Quantile Sensitivity of SOFR-RRP Spread to Reserves and Dealer Positions Notes: The panels present the regression results from Equation 1.}
\end{figure}
\subsection{AI、技术与生产率}
\textbf{用于判断 AI 投资周期、生产率扩散和产业约束是否正在改变资产定价。}
\begin{itemize}
  \item 人工智能正在为全球宏观环境节省巨额劳动成本。IMF最新工作论文基于Anthropic宏观环境指数五轮数据（2025年1月至2026年2月）估算，当前AI使用带来的时间节省价值约为每年2.7万亿美元，相当于样本国家GDP总和的3.4\%。但这一数字掩盖了一个结构性不对称：在低收入国家...
\end{itemize}
\begin{figure}[htbp]
\centering
\includegraphics[width=0.88\linewidth]{figures/fig_100.jpg}
\caption{Figure 2 - Figure 2: Aggregate AI gains: scale, trend, and incidence by income. Panel (a) (R1 to R5, Jan 2025 to Feb 2026): the blue line (left axis) is the wage index, the average US occupation-level wage weighted by global AI conversation}
\end{figure}
\section{结语}
后续需验证的数据与叙事节点包括：美国6月CPI数据对近端利率的非对称风险释放；中国7月政治局会议对财政与消费政策的定调；AI基础设施电力瓶颈的缓解信号（如燃气轮机订单与电网投资）；Netflix二季度财报对用户参与度与广告收入的验证；欧洲对中国电动车关税政策的进一步动向。
\newpage
\section{覆盖概览}
投行/券商 40 篇；战略咨询 2 篇；智库/国际机构 23 篇。\par
投行覆盖：GS 12篇 / Citi 10篇 / NOM 7篇 / Bernstein 5篇 / JPM 4篇 / DB 1篇 / UBS 1篇\par
\section{Disclaimer}
{\scriptsize\color{refgray}
This community is a paid learning and information-sharing community created for general educational purposes only. The membership fee is charged solely for access to community discussions, educational content, organization, and operational costs. It is not an advisory fee, management fee, performance fee, brokerage fee, or compensation for personalized financial, investment, legal, tax, or professional advice.\par
I participate and share information solely as an individual and for educational discussion among community members. I am not acting as a financial adviser, investment adviser, broker, dealer, portfolio manager, legal adviser, tax adviser, accountant, fiduciary, or any other licensed professional adviser. Nothing shared in this community constitutes, and should not be interpreted as, financial advice, investment advice, legal advice, tax advice, a recommendation, solicitation, offer, endorsement, instruction, or guarantee to buy, sell, hold, short, trade, or invest in any security, cryptocurrency, fund, derivative, strategy, or other financial product.\par
All content shared in this community, including messages, discussions, links, files, charts, examples, market commentary, watchlists, AI-generated or AI-polished summaries, and third-party materials, is general, impersonal, and educational in nature. It does not take into account any individual's financial situation, investment objectives, risk tolerance, investment horizon, liquidity needs, tax status, legal restrictions, or suitability requirements. No content should be relied upon as suitable for any specific person, account, portfolio, or financial decision.\par
Information shared in this community may be incomplete, inaccurate, outdated, speculative, AI-generated, AI-edited, or based on third-party internet sources that have not been independently verified. I make no representation or warranty as to the accuracy, completeness, timeliness, reliability, or suitability of any information shared.\par
Financial markets involve substantial risk, including the possible loss of principal. Past performance, historical data, hypothetical examples, simulated results, backtests, personal opinions, or educational examples do not guarantee future results. Each member is solely responsible for conducting their own research, exercising independent judgment, and making their own decisions. Before making any financial, investment, legal, or tax decision, members should consult qualified and licensed professionals.\par
Participation in this community, payment of any membership fee, or communication with me or other members does not create any adviser-client, fiduciary, professional, contractual, agency, partnership, employment, or other advisory relationship. By joining or participating in this community, you acknowledge and agree that you are solely responsible for your own decisions, actions, transactions, profits, losses, risks, and consequences, and that I shall not be liable for any direct, indirect, incidental, consequential, financial, legal, tax, or other loss or damage arising from or related to any content, discussion, information, or material shared in this community.\par
}
\newpage
\begin{center}
{\Large 更多详情报告kcdesk.com}\par\vspace{0.8cm}
\includegraphics[width=0.58\linewidth]{\detokenize{../../prompts/zsxq_img.jpg}}
\end{center}
\end{document}
